\section{Resolución del problema}

La evaluación muestra que la regla hace lo que se diseñó para hacer. Aquí esa regla deja de ser solo un cálculo y se vuelve el recorrido de quien compra. Los datos, las decisiones de la exploración y el orden llegan a la persona en una aplicación. No es un proyecto aparte del análisis.

El problema original era comparar alimentos empacados cuando la información está repartida, sin tratar un silencio como si fuera un dato. El recorrido que lo resuelve es el que la estrategia ya nombró: buscar, entender, comparar alternativas, decidir y cerrar la compra.

La búsqueda es por nombre o por código de barras, sobre todo el catálogo operativo, antes de calcular el orden. Así se puede llegar a un producto aunque después no entre a la lista.

En la ficha se entiende el resultado. Se ven nutrición, procesamiento y etiquetas en lenguaje cotidiano, el estado de alergia y de dieta, y el precio con su procedencia. Si el precio no fue observado, se marca como demostración. No es un precio de tienda y no ordena. La banda dice si el producto entra en la comparación, si la información no alcanza, si la alergia o la dieta no se pueden verificar, o si queda excluido. El detalle de cada nutriente existe en el cálculo. La ficha muestra las tres señales y las razones de alergia y dieta. No reescribe, nutriente por nutriente, el texto que ese cálculo ya produjo.

Desde la ficha se pasa a otras opciones de la misma categoría. Esa sustitución es la comparación para la que se construyó el grupo: productos parecidos, no el catálogo entero. La persona conserva el producto, lo cambia o lo deja fuera. La aplicación no cierra esa decisión.

Si elige una opción, la lleva al carrito. El carrito resume el conjunto y lo reparte en los grupos del Plato del Bien Comer: verduras y frutas, cereales, leguminosas y alimentos de origen animal, y no clasificado \citep{nom0432013}. Si la categoría de un producto no se puede ubicar, queda en no clasificado. El resumen de partida es el promedio por 100\,g. El promedio que pesa los gramos de cada producto solo se usa cuando todos tienen cantidad. El resumen dice qué método usó y qué tanta información había. Las alertas avisan porciones y sellos cuando el producto los declara. El resumen no emite un juicio médico.

También hay una vista de todo el catálogo, sin partir de una categoría. Sirve para explorar. No sustituye la comparación entre productos parecidos.

La aplicación, en este punto, es la forma en que el análisis llega a la compra. Deja un orden, una banda y alternativas del mismo tipo, construidos solo con la información que el registro sí trae.
